\chapter{Visual Field Test (Perimetry)}

\section*{Overview}
A visual field test measures how much a person can see to the sides, above, and below while looking at a fixed central target. It detects areas of reduced or missing vision that may not be noticed in daily life.

\section*{Why It Is Done}
It is commonly used to diagnose and monitor glaucoma, optic-nerve disease, retinal disorders, and neurological conditions affecting the visual pathways. It may also document changes caused by medicines or other eye disease.

\section*{How It Is Performed}
One eye is tested at a time while the patient looks at a central point inside a bowl-shaped instrument. Small lights appear in different locations, and the patient presses a button whenever a light is seen. The test is repeated or adjusted if attention, fatigue, or learning affects reliability.

\section*{Preparation and Results}
Wear usual glasses unless instructed otherwise and bring information about eye conditions and medicines. Results are shown as a map of sensitivity and are interpreted with the eye examination, intraocular pressure, and previous tests. A reliable test may require good attention and steady fixation.

\section*{Risks and Limitations}
The test is noninvasive and painless. Temporary eye fatigue or frustration can occur. Poor concentration, blurred vision, uncorrected refractive error, or eyelid obstruction can produce an apparently abnormal result.

\section*{Contraindications and Precautions}
There are usually no major medical contraindications. Testing may need modification for severe fatigue, inability to understand instructions, uncorrected vision problems, eyelid obstruction, or inability to maintain fixation. The result is interpreted with the eye examination and reliability indicators.

\section*{When to Seek Medical Care}
Follow the procedure team's specific instructions. Seek urgent medical assessment for severe or worsening pain, uncontrolled bleeding, fever or signs of infection, trouble breathing, chest pain, fainting, new weakness or confusion, inability to urinate, or any rapidly worsening symptom. The appropriate warning signs depend on the procedure and the patient's medical condition.

\section*{References}
\begin{enumerate}
  \item National Eye Institute. Visual field test. 2025.
  \item Current Medical Diagnosis and Treatment. McGraw-Hill; 2025.
\end{enumerate}
\clearpage
